\begin{abstractpage}
\vspace*{-3cm}

\begin{abstract}{romanian}
Domeniul vederii artificiale este reprezentat de extragerea informațiilor din imagini, folosind contextul vizual pentru a sprijini alte sisteme software automate. Astfel, unul dintre domeniile promițătoare este reprezentat de competițiile sportive deoarece regulile stricte și necesitatea deciziilor rapide constituie atât un avantaj, cât și o provocare.

Proiectul de diplomă prezentat își propune să introducă un sistem automat de analiză a jocurilor de tenis de masă, FMI-TT, care orchestrează mai multe arhitecturi de rețele neuronale pentru detecția elementelor: GDINO pentru propuneri de regiuni, SAM pentru segmentare semantică, YOLO pentru detectarea elementelor dinamice, dar și tehnici clasice de vedere artificială, precum filtrarea HSV, dilatarea și eroziunea. În același timp, vom adapta o bază de date în care sunt prezente cadre din diferite jocuri de tenis de masă adnotate pentru antrenarea și validarea sistemului propus.

Sistemul automat va detecta mingea, paletele și masa de joc, extrăgând informații relevante pentru a identifica evenimentele-cheie, cum ar fi loviturile de retur, contactul mingii cu masa și traiectoria acesteia.

Pentru a cuantifica rezultatele obținute, vom realiza o comparație între datele experimentale obținute și cele prezentate în TTNet, o lucrare științifică realizată de institutul OSAI, care propune un sistem de analiză automată a jocurilor de tenis de masă, dar utilizând o abordare diferită bazată pe analiza întregului joc de tenis de masă cu rețele neuronale.

\end{abstract}

\begin{abstract}{english}
The field of computer vision involves extracting information from images by using visual context to support other automated software systems. One of the most promising areas of application is sports competitions, as their strict rules and need for rapid decision-making represent both an advantage and a challenge.

This bachelor’s thesis aims to introduce an automated system for analyzing table tennis matches, named FMI-TT, which orchestrates multiple neural network architectures for element detection: GDINO for region proposals, SAM for semantic segmentation, YOLO for detecting dynamic elements, as well as classical computer vision techniques such as HSV filtering, dilation, and erosion. At the same time, we will adapt a database containing frames from various annotated table tennis matches for training and validating the proposed system.

The automated system will detect the ball, paddles, and table, extracting relevant information to identify key events such as returns, ball-table contact, and its trajectory.

To evaluate the obtained results, we will perform a comparison between our experimental data and the results presented in TTNet, a scientific work developed by the OSAI institute, which proposes an automatic analysis system for table tennis matches, but using a different approach based on whole-match analysis through neural networks.


\end{abstract}

\end{abstractpage}
